% !TEX program = xelatex
\documentclass[12pt,a4paper]{article}

\usepackage{fontspec}
\setmainfont{Arial} % FAPESP APR: Calibri or Arial font, size 12
\usepackage[english]{babel}
\usepackage{microtype}
\usepackage[a4paper,left=3cm,right=1.5cm,top=2.5cm,bottom=2.5cm]{geometry} % FAPESP APR: 3cm left / 1.5cm right margins
\usepackage{setspace}
\onehalfspacing % FAPESP APR: 1.5 line spacing
\usepackage{amsmath,amssymb}
\usepackage{booktabs}
\usepackage{longtable}
\usepackage{array}
\usepackage{enumitem}
\usepackage{xcolor}
\usepackage{titlesec}
\usepackage[numbers,sort&compress]{natbib}
\usepackage[hidelinks,breaklinks]{hyperref}

\setlist{itemsep=2pt,topsep=4pt}
\titleformat{\section}{\normalfont\large\bfseries}{\thesection}{0.75em}{}
\titleformat{\subsection}{\normalfont\normalsize\bfseries}{\thesubsection}{0.75em}{}

\newcommand{\pamr}{P\nobreakdash-AMR}

\title{\textbf{Precision-AMR: Spatially Adaptive Floating-Point Precision
for PDE-Based Reverse Time Migration}\\[0.5em]
\large\normalfont\itshape A new axis of numerical adaptivity for adjoint
wave-equation solvers, developed on dedicated NVIDIA hardware and validated
across heterogeneous NVIDIA and AMD GPU architectures on the Santos Dumont
supercomputer (LNCC)}

\author{}
\date{}

\begin{document}

\maketitle
\thispagestyle{empty}

\begin{flushleft}
\textbf{Principal Investigator:} Prof.\ Sandro Rigo\\
\textbf{Title/Role:} Associate Professor, Institute of Computing\\
\textbf{University:} Universidade Estadual de Campinas (UNICAMP)\\[0.4em]
\textbf{Collaborator 1:} Thiago Maltempi --- Graduate Researcher (PhD
Candidate), Institute of Computing, Universidade Estadual de Campinas\\
\textbf{Collaborator 2:} Prof.\ Herv\'e Yviquel --- Assistant Professor,
Institute of Computing, Universidade Estadual de Campinas
\end{flushleft}

\section{Problem Statement}

GPU hardware is increasingly shifting toward low precision in favor of AI
workloads (FP8/BF16), exposing new memory hierarchies, tensor cores, and
mixed-precision math libraries. At the same time, scientific simulations
based on partial differential equations (PDEs) continue to require FP64/FP32
precision to remain trustworthy. Reverse Time Migration (RTM), a seismic
imaging method that solves the wave equation by finite differences over
thousands of time steps, is a canonical example of this tension:
geophysicists use the migrated image to guide expensive drilling decisions,
so numerical fidelity is essential. RTM is also memory-intensive, requiring
hundreds of gigabytes of wavefield state and relying on checkpointing to
trade memory against recomputation in its forward/backward adjoint
structure.

Our prior work tackled this structure's main bottleneck, namely host--GPU
checkpoint transfer over PCIe, by combining the concept of prefetching with
data compression. This idea was materialized as the GPUZIP
library~\cite{maltempi2025checkpointing,gpuzip}, which compresses checkpoints
on the GPU and prefetches them proactively, reducing host--GPU transfer
volume by up to 80$\times$ and achieving end-to-end speedups of up to
5.1$\times$ on V100, while preserving image quality (SSIM $\approx$ 1,
near-zero relative error). Separately, reduced-precision arithmetic has
already been shown to be viable for seismic modeling, RTM, and FWI:
Fabien-Ouellet~\cite{fabienouellet2020half} found that a stable
half-precision (FP16) implementation of the elastic wave equation is
achievable through appropriate scaling, that the resulting seismogram error
is a negligible fraction of total seismic energy and largely incoherent with
seismic phases, and that this noise does not adversely affect FWI or RTM
image quality. However, that reduction is applied uniformly across the
entire computational domain. Gao and Harms~\cite{gao2023half} identified a
specific hazard of that approach: naive half-precision time stepping
accumulates round-off error through a disguised recursive summation in the
update rule, degrading solution quality, with compensated (Kahan) summation
as an effective remedy.

At the same time, adaptive mesh refinement (AMR) has been explored for
seismic wave simulation and, more recently, explicitly for FWI and RTM
workflows, to avoid the waste inherent in a uniform grid: dispersion and
dissipation errors require grid spacing proportional to the shortest
wavelength, so a uniform mesh is forced to use the finest spacing
everywhere, even where the local velocity field would tolerate a much
coarser one. AMR addresses this by refining $\Delta x$ locally. No
comparable, physically motivated scheme exists for floating-point precision,
however: existing mixed-precision work in seismic imaging is either global
and uniform (as above), or adaptive only at the level of an individual
linear-algebra operation (mixed-precision GEMM, iterative refinement,
Ozaki-scheme emulation), not tied to a spatial decomposition of the
wavefield itself. None of this prior work varies floating-point precision
according to physically motivated spatial criteria. This is therefore a
third axis of adaptivity that remains largely unexplored.

This project proposes to close that gap with Precision-AMR (\pamr{}): a
spatially block-structured scheme, an adaptive-precision analogue of AMR,
that keeps mesh resolution fixed and instead varies the floating-point
format assigned to each block (tile) of the domain, based on physically
motivated error-tolerance criteria (proximity to absorbing boundaries,
wavefield energy density, adjoint-sensitivity magnitude, and a
time-dependent band trailing the propagating wavefront) that predict whether
the round-off error introduced in that block can reach the imaging
condition or the FWI gradient. This is a structurally different problem
from mesh-based AMR. In AMR, the characteristic artifact is a localized
reflection at the coarse-fine grid interface, because refinement changes
the discretization itself. In \pamr{}, no discretization boundary is
introduced; instead, a spatially varying level of round-off noise is
injected into a hyperbolic PDE, and that noise propagates together with the
wave. Tagging criteria therefore cannot rely solely on local smoothness, as
AMR refinement indicators do; they must instead approximate where an error,
once injected, will decay, be absorbed, or otherwise fail to reach the
receivers or the sensitivity kernel before the simulation ends, and control
round-off accumulation over the thousands of time steps typical of RTM/FWI.

Because RTM/FWI already require a block/tile decomposition of the domain
for checkpointing, \pamr{} reuses that same decomposition as its tagging
granularity, unifying the precision budget and the lossy-compression budget
of differential checkpointing into a single, jointly tuned per-tile accuracy
target. Mixed-precision GEMM, cuSOLVER iterative refinement, and
Ozaki-scheme FP64 emulation remain essential as supporting techniques, but
are now applied selectively, only where \pamr{}'s criteria indicate they are
needed. This is the natural way to exploit GPUs geared primarily toward AI,
such as the RTX 5090, which offer negligible native FP64 throughput.

Because the tagging criteria are motivated by the physics of wave
propagation and adjoint sensitivity, rather than by any particular vendor's
arithmetic units, we treat their portability across architectures as an
explicit question, not an assumption. Our research group has access to the
Santos Dumont supercomputer (LNCC/SINAPAD), which includes AMD Instinct
MI300A accelerators alongside several NVIDIA generations (Volta V100, Hopper
H100, Grace-Hopper GH200). This lets us test that question directly, rather
than merely assert it, and it is what gives Precision-AMR its potential
contribution to the field: a numerical adaptivity axis that is new and
portable across vendors, not merely a one-off optimization for a single GPU
architecture.

\section{Expected Results}

The project's principal contribution is the Precision-AMR (\pamr{}) method
itself, which explores a new axis of numerical adaptivity that varies
floating-point precision spatially, according to physically motivated
criteria (proximity to the PML, wavefield energy density, adjoint-sensitivity
magnitude, and a wavefront-trailing band), proposed and validated
specifically for adjoint PDE solvers such as RTM/FWI. This contribution
includes the formulation of the four tagging criteria, the characterization
of the ``precision-interface'' artifact, analogous to the coarse-fine
interface reflections of classical AMR, and the demonstration that the
precision and lossy-compression budgets can be unified into a single
per-tile error target.

As secondary contributions, the project will produce: (i) a
multi-architecture evaluation of P-AMR-guided mixed-precision usage,
comparing NVIDIA GPUs (V100, H100, GH200) and the AMD Instinct MI300A APU,
characterizing how much of the achieved speedup and accuracy trade-off is
vendor-independent; and (ii) the implementation of P-AMR as an
\textbf{independent library}, with an interface that allows its integration
into other PDE-based simulation applications, beyond the specific scope of
RTM/FWI. This library may be adopted standalone or combined with GPUZIP,
merging P-AMR's spatial precision zoning with GPUZIP's compressed
checkpointing into a single workflow, without requiring the host application
to adopt the entire GPUZIP stack just to benefit from P-AMR. It will be made
available as open source, under the MIT license, on GitHub.

The central scientific result will be submitted to a high-impact
high-performance-computing event or journal, for example, Supercomputing
(SC), IPDPS, Euro-Par, or SBAC-PAD among the events, or IJHPCA or a
Transactions-type journal from IEEE or ACM among the journals, accompanied
by a public reproducibility package (datasets, Nsight profiles, quality
metrics), made available via UNICAMP's REDU.

\section{Scientific and Technological Challenges and the Means and Methods
to Overcome Them}

The feasibility of adaptive mesh refinement (AMR) for reducing the
computational cost of seismic simulation is already well established in
the literature. van Driel et al.~\cite{vandriel2020wavefield1} and
Thrastarson et al.~\cite{thrastarson2020wavefield2} show, respectively,
that wavefield-adapted meshes accelerate forward/adjoint modelling and FWI
by up to an order of magnitude when medium properties vary smoothly. Yang
et al.~\cite{yang2024elasticrtm} apply a variable-spacing mesh specifically
to elastic RTM, preserving multicomponent imaging accuracy while reducing
computational cost and storage. Davis and LeVeque~\cite{davis2015adjoint}
further show that refinement indicators guided by the adjoint computation,
and not just by local smoothness of the field, produce more efficient mesh
refinements in hyperbolic wave-propagation problems, the same logic that
motivates \pamr{}'s sensitivity (FWI) criterion. None of these works,
however, varies \emph{floating-point precision} instead of \emph{mesh
resolution}. Nor does any of them characterize empirically how the
numerical behavior of an RTM/FWI simulation, that is, the dynamic range of
floating-point values actually exercised, distributes across space and
time. This double gap, the absence of a physically motivated
precision-zoning scheme and the absence of the empirical characterization
needed to ground it before designing the criteria, defines this project's
scientific and technological challenges, organized in the six phases
below, over a 36-month programme.

\subsection{Phase 1 --- Infrastructure and Baseline (M1--M4)}

\textbf{Challenge:} before any precision change is attempted, it is
necessary to faithfully reproduce the GPUZIP/Awave-3D baselines under the
HPC SDK and ensure that the block decomposition used for checkpointing can
also carry a per-block precision tag, an engineering prerequisite that
every subsequent phase depends on.

\textbf{Method:} port Awave-3D and GPUZIP to the HPC SDK (nvc++),
re-integrate nvCOMP/Bitcomp, and reproduce the V100 baselines on Santos
Dumont's V100 partition, matching the original GPUZIP hardware.

\begin{itemize}
  \item Establish a block/tile decomposition of the domain, extending the
        existing checkpoint-block structure so that it can also carry a
        precision tag per block.
  \item Profile every kernel with Nsight Systems and Nsight Compute, on
        both the RTX 5090 and Santos Dumont, confirming that FDM stencils
        are memory-bandwidth-bound at the tile granularity before any
        precision change is attempted.
\end{itemize}

\subsection{Phase 2 --- Empirical Characterization of Mixed-Precision
Behavior (M4--M10)}

\textbf{Challenge:} the four tagging criteria proposed for Phase 3 are
physically motivated, but, unlike the adjoint-guided mesh refinement of
Davis and LeVeque~\cite{davis2015adjoint}, no work in the literature
empirically characterizes how the dynamic range of floating-point values
(exponents and magnitudes actually used) behaves spatiotemporally in an
RTM/FWI simulation, nor investigates whether that behavior can be
anticipated directly from the input velocity model, before the wave even
propagates. Without that characterization, the tagging criteria remain
plausible hypotheses that have not yet been verified against actual
numerical behavior.

\textbf{Method:} instrument the solver to record, from FP64 reference
runs, dynamic-range statistics of the wavefield per block and per time
step: exponent histogram, minimum/maximum magnitude, variance, across the
six datasets described in the Evaluation section.

\begin{itemize}
  \item Correlate these statistics with attributes derived from the input
        velocity model (local impedance contrast, depth, distance to
        reflectors, local wavelength), seeking which attributes are
        predictive of where the dynamic range, and therefore the
        tolerance to round-off error, is structurally lower.
  \item From that correlation, build an \emph{a priori} precision-zoning
        estimator, computable directly from the velocity model before any
        time step is run. This estimator is complementary to the four
        reactive criteria of Phase 3, which depend on the wavefield
        computed during the simulation itself.
  \item Use this characterization to validate, or refute, the physical
        assumptions behind each of the four proposed criteria, before
        investing in their implementation and fine-tuning in Phase 3.
\end{itemize}

\subsection{Phase 3 --- Precision-AMR: Criteria, Tagging, and Validation
(M10--M20)}

\textbf{Challenge:} unlike mesh-based AMR, where the characteristic
artifact is a localized reflection at the coarse-fine interface, round-off
noise induced by precision reduction propagates together with the wave
itself, so tagging criteria cannot rely solely on local smoothness of the
field, as classical refinement indicators do. They must anticipate where
an error, once injected, will fail to reach the receivers or the FWI
gradient before the simulation ends. A second, technological challenge
compounds this: reduced precision and lossy compression, both already used
separately in GPUZIP, may interact non-trivially when applied to the same
block, and that interaction needs to be isolated before defining a joint
error budget.

\textbf{Method:} this is the core scientific phase, run at scale on Santos
Dumont's H100 partition. Building on the characterization from Phase 2, we
design, implement, and validate four candidate tagging criteria for
spatial precision zoning:

\begin{itemize}
  \item \textbf{Boundary criterion:} tiles inside or adjacent to the
        PML/absorbing layer are tagged for reduced precision by default,
        since their contribution is already being damped.
  \item \textbf{Energy criterion:} tiles are tagged by a threshold on
        local wavefield energy density, so that regions of low amplitude
        (before wavefront arrival, or after geometrical spreading has
        attenuated the signal) tolerate lower precision.
  \item \textbf{Sensitivity criterion (FWI):} tiles are tagged using the
        magnitude of the local adjoint-sensitivity kernel, reusing the
        gradient computation itself as the tagging signal, in direct
        analogy to the adjoint-guided refinement of Davis and
        LeVeque~\cite{davis2015adjoint} and to the Berger--Oliger
        truncation-error estimator in classical AMR.
  \item \textbf{Wavefront-trailing criterion:} a time-dependent band
        immediately behind the propagating wavefront is tagged for
        reduced precision, while the leading edge, where phase accuracy
        governs the imaging condition, retains full precision.
\end{itemize}

For each criterion, we implement compensated (Kahan) summation in the
low-precision tiles' time-stepping update to control round-off
accumulation over the full simulation, following the mechanism identified
by Gao and Harms~\cite{gao2023half}. We explicitly test for a
``precision-interface'' artifact, analogous to coarse-fine interface
reflections in mesh-based AMR, via residual analysis at the boundaries
between tiles of differing precisions.

To isolate the effects of mixed precision from the effects of lossy
compression before unifying them into a single per-tile error budget, each
candidate tagging region is evaluated under four conditions, always
compared against the FP64 reference by PSNR, SSIM, and mean relative
error, on both the raw wavefield and the final migrated image:

\begin{enumerate}
  \item full precision, no compression (control);
  \item full precision, with differential compression (isolates the
        compression effect);
  \item reduced precision, no compression (isolates the precision
        effect);
  \item reduced precision, with differential compression (the condition
        actually deployed in production).
\end{enumerate}

Comparing conditions (ii)+(iii) against (iv) indicates whether the two
effects are approximately additive, whether one masks the other, or
whether there is amplification through interaction. Only then does each
tile jointly receive a floating-point format and a checkpoint compression
policy (full snapshot versus compressed delta), as a single, jointly
tuned error budget. This phase subsumes differential checkpointing, which
becomes the storage mechanism for \pamr{}-tagged tiles rather than a
separate deliverable.

\subsection{Phase 4 --- Precision-AMR without Native FP64 (M19--M27)}

\textbf{Challenge:} GPUs geared primarily toward AI, such as the RTX 5090,
offer negligible native FP64 throughput. Emulating FP64 across the entire
domain via the Ozaki scheme~\cite{uchino2025ozaki} is technically viable
but expensive enough to make indiscriminate use impractical, which only
reinforces the need to apply it selectively.

\textbf{Method:} the tagging criteria validated in Phase 3 are applied on
the RTX 5090. Per tile, the pipeline selects between native low-precision
tensor-core execution, cuBLAS mixed-precision GEMM (FP32 compute with
FP64 accumulation) or cuSOLVER iterative refinement for dense sub-kernels,
and Ozaki-scheme FP64 emulation for tiles flagged as accuracy-critical.

\begin{itemize}
  \item This directly operationalizes the motivation for exploiting
        AI-first GPUs that lack native FP64: expensive emulation is
        concentrated only where \pamr{}'s criteria indicate it is needed.
  \item We quantify the speedup/accuracy trade-off relative to a
        uniform-Ozaki baseline (emulating FP64 everywhere) to determine
        how much selective tiling actually saves.
  \item All errors are bounded jointly (compression $\times$ precision
        interaction, following the Phase 3 methodology) against the FP64
        baseline using PSNR, SSIM, and mean relative error, across the
        same six datasets and three checkpointing algorithms (Revolve,
        zCut, Uniform) used throughout the project.
\end{itemize}

\subsection{Phase 5 --- Cross-Architecture Validation (AMD Instinct
MI300A) (M26--M31)}

\textbf{Challenge:} the tagging criteria are motivated by the physics of
wave propagation and adjoint sensitivity, not by any particular vendor's
arithmetic units. That portability, however, is a claim that needs to be
tested, not merely assumed. The concrete obstacle is the compression
stack: GPUZIP today compresses checkpoints with Bitcomp, which is
NVIDIA-proprietary and has no HIP port, so there is no checkpoint
compression support on AMD GPUs. Generic \emph{lossless} compression (LZ4,
Snappy, Cascaded) does not solve the problem: Künas et
al.~\cite{kunas2026amdcompression} ported these algorithms from CUDA to
HIP and measured, on AMD CDNA and RDNA GPUs (including the MI300X, of the
same CDNA3 family as the MI300A), a compression ratio of only 1.03 to
1.04$\times$ for floating-point seismic data, regardless of architecture.
This is why GPUZIP uses bounded-error lossy compression, and it is this
component that must be rebuilt for AMD.

\textbf{Method:} Phase 5 begins by incorporating into GPUZIP a compression
library with native AMD support, replacing the Bitcomp/nvCOMP path on the
CDNA3 architecture. The options under consideration are:

\begin{itemize}
  \item \textbf{hipCOMP-core}~\cite{hipcompcore}: a HIP port of the same
        algorithms as nvCOMP (LZ4, Snappy, Cascaded), with HIP/AMD and
        HIP/CUDA backends. It is the direct analogue of nvCOMP, but still
        \emph{early-access} and not optimized for AMD GPUs, and it covers
        only the \emph{lossless} case.
  \item \textbf{HPDR}~\cite{chen2025hpdr}: a portable framework with
        native adapters for OpenMP, CUDA, and HIP, implementing pipelines
        of MGARD (multilevel \emph{lossy} compression with error bounded
        in quantities of interest), ZFP (fixed-rate compression for
        multidimensional arrays), and Huffman. It was evaluated on up to
        1,024 nodes of Frontier (AMD Instinct MI250X GPUs), reaching data
        reduction throughput of up to 103\,TB/s. It is the most mature
        functional equivalent of Bitcomp on Instinct hardware, covering
        the bounded-error \emph{lossy} compression role.
\end{itemize}

The final choice between these alternatives, or another that emerges, is
made at the start of this phase, according to the maturity of AMD
compression libraries at that time. The current guidance is to target
HPDR with an MGARD or ZFP pipeline as the Bitcomp replacement. With the
compression stack ported, the validation proper proceeds:

\begin{itemize}
  \item We port the same tagging and tiling logic to Santos Dumont's AMD
        Instinct MI300A nodes via ROCm/HIP (rocBLAS, hipBLAS, and AMD's
        own reduced-precision matrix paths), without assuming feature
        parity with the CUDA-specific path.
  \item We test whether the Phase 3 criteria and the precision $\times$
        compression factorial methodology hold on a CDNA3 architecture,
        re-running the same evaluation protocol used on the NVIDIA GPUs,
        now with the AMD compressor in place of Bitcomp.
  \item We test whether the \emph{a priori} estimator built in Phase 2,
        derived from the velocity model and architecture-independent,
        transfers directly, without retuning, to the MI300A.
  \item We compare the speedup/accuracy trade-off achieved on the MI300A
        against the NVIDIA baselines, to characterize how much of the
        gain is architecture-independent and how much is specific to the
        CUDA/Ozaki/nvCOMP stack.
\end{itemize}

\subsection{Phase 6 --- Spectral Operators, Integration, and Release
(M31--M36)}

\textbf{Challenge:} consolidating the supporting techniques (differential
checkpointing, mixed-precision GEMM, Ozaki-scheme emulation, spectral
operators) and the Precision-AMR module into a single, coherent, reusable
software artifact, the P-AMR library described in Expected Results,
rather than leaving them as isolated prototypes from each phase.

\textbf{Method:}

\begin{itemize}
  \item Benchmark cuFFT-based pseudo-spectral derivative operators as an
        alternative update rule for selected, precision-tagged tiles, as
        a secondary technique complementing the core \pamr{} contribution.
  \item Integrate the Precision-AMR module, already as an independent
        library, differential checkpointing, and mixed-precision paths
        into a unified GPUZIP v3.0 release.
  \item Multi-node scaling evaluation on Santos Dumont ($\leq$8 GPUs).
  \item Final analysis, write-up, and submission, with Precision-AMR as
        the central scientific claim and differential checkpointing,
        mixed-precision linear algebra, and spectral operators presented
        as supporting techniques.
\end{itemize}

\subsection{Computational Platforms}

\paragraph{HPC SDK} (nvc++/nvfortran, CUDA-X math libraries, Nsight
Systems/Compute). Experience: team uses CUDA and Nsight Systems routinely.
Adoption: migrate the Awave-3D build from clang to nvc++ and validate
numerical parity against existing baselines. Nsight Compute provides
kernel-level roofline analysis, now applied per spatial tile, to identify
which tiles are memory-bandwidth-bound versus compute-bound and to
objectively scope where a lower-precision tile actually pays off.

\paragraph{RTX 5090} (dedicated hardware, requested, 32\,GB GDDR7 plus a
server with 64\,GB of RAM). This GPU is also based on the Blackwell
architecture and shares with same-generation professional cards the same
memory bandwidth (1,792\,GB/s) and the same 5th-generation tensor cores,
with FP8/FP4 support and negligible native FP64 throughput. For a
memory-bandwidth-bound finite-difference solver like ours, it is memory
bandwidth, not raw CUDA-core count, that determines performance. This
makes the RTX 5090 a representative development platform for \pamr{}, at
a fraction of the cost of an equivalent professional card and within the
per-item equipment ceiling set by FAPESP's Regular Research Grant rules
for this type of request. The absence of fast native FP64 is not a new
disadvantage: it is exactly the precision regime that \pamr{} already
assumes and exploits, via Ozaki-scheme emulation and mixed-precision
GEMM, so the scientific question this project answers, precisely which
tiles need emulated FP64 and which do not, remains the same. Its 32\,GB
of GDDR7 (without ECC) support kernel iteration and the tuning of tagging
criteria on representative subsets of the domain; as dedicated,
queue-free hardware, it enables fast iterative debugging that a shared
batch system is not well suited for. FP64 ground truth and production
validation continue to be generated on Santos Dumont, with ECC-protected
memory.

\paragraph{Santos Dumont} (LNCC/SINAPAD, Tier-0 national supercomputer,
Petr\'opolis-RJ). Provides the large-scale validation environment. Its
BullSequana X1000 partition includes 94 nodes with 4$\times$ NVIDIA Volta
V100 GPUs each, matching the hardware used for the original GPUZIP
baselines, plus a dedicated AI node with 8$\times$ V100 16GB over NVLink.
Its newer BullSequana XH3000 partition adds 62 nodes with 4$\times$ NVIDIA
Hopper H100 SXM 80GB over NVLink, 36 NVIDIA Grace-Hopper GH200 superchip
nodes, and 18 nodes with 2$\times$ AMD Instinct MI300A APUs (CDNA3 GPU
cores, 128\,GB HBM3 each). This heterogeneity lets us generate FP64 ground
truth and run the full \pamr{} criteria sweep at scale on NVIDIA hardware,
and separately test whether the tagging methodology and block-tiling
approach generalize to a non-NVIDIA architecture.

\paragraph{nvCOMP / Bitcomp} (GPU-side lossy compressor). Experience:
integrated and validated in GPUZIP v2.0, where Bitcomp delivered the best
accuracy-compression balance on V100. Adoption here: the storage mechanism
for differential checkpointing, now indexed by the same tile identifiers
used for precision tagging, so each tile carries a single, jointly-tuned
precision-and-compression error budget.

\paragraph{cuBLAS (mixed-precision extensions) and cuSOLVER (iterative
refinement).} Both ship inside the HPC SDK. Experience: team has used
BLAS/LAPACK routines in prior HPC work. Adoption: applied only within
tiles that \pamr{} flags as containing dense sub-kernels needing accuracy
beyond native low precision (FP32 compute with FP64 accumulation, and
cuSOLVER iterative refinement).

\paragraph{Ozaki-scheme FP64 emulation} via low-precision tensor cores.
Adoption: reserved for the subset of tiles that \pamr{}'s sensitivity
criteria mark as accuracy-critical (near the wavefront, high
adjoint-sensitivity magnitude), rather than applied uniformly. This
directly tests whether selective emulation is cheaper than emulating
everywhere while preserving image fidelity. Recent work has demonstrated
up to 2.3$\times$ speedup over native FP64 DGEMM on Blackwell
GPUs~\cite{uchino2025ozaki}, measured on the professional RTX PRO 6000
card; we expect comparable results on the RTX 5090, since both share the
same tensor-core generation.

\paragraph{cuFFT.} Experience: used in prior spectral operator work.
Adoption: benchmarked as an alternative to FDM stencils for wavefield
updates within selected, precision-tagged tiles, characterising the
memory-bandwidth versus compute trade-off relative to the standard
finite-difference path.

\paragraph{AMD Instinct MI300A} (CDNA3, Santos Dumont, exploratory). As an
additional validation activity, we assess how much of the \pamr{}
framework, that is, the block/tile decomposition, the four tagging
criteria, and the joint precision-compression error budget, transfers to
AMD's ROCm/HIP ecosystem (rocBLAS, hipBLAS, and AMD's own
reduced-precision matrix paths), without assuming feature parity with the
CUDA-specific Ozaki-scheme emulation or nvCOMP/Bitcomp compression path.
This last gap is concrete: there is no HIP port of Bitcomp,
hipCOMP-core~\cite{hipcompcore} is still \emph{early-access}, and generic
\emph{lossless} compression yields only 1.03 to 1.04$\times$ on
floating-point seismic data~\cite{kunas2026amdcompression}, which
motivates using a portable bounded-error \emph{lossy} compressor such as
HPDR~\cite{chen2025hpdr} in Phase 5. The goal is to determine whether
Precision-AMR is a portable methodology or an NVIDIA-specific
optimization.

\section{Execution Timeline}

The project runs for 36 months, organized in the six phases described in
the previous section. Completion is planned for month 36, with paper
submission and the GPUZIP v3.0 release. Each phase ends with a verifiable
deliverable that serves as a progress milestone; the v3.0-alpha release at
the end of Phase 3 (month 20) is the main intermediate milestone.

\begin{longtable}{@{}p{0.07\textwidth}p{0.08\textwidth}p{0.44\textwidth}p{0.30\textwidth}@{}}
\toprule
\textbf{Phase} & \textbf{Months} & \textbf{Main activities} &
\textbf{Deliverable at the end of the phase}\\
\midrule
\endfirsthead
\toprule
\textbf{Phase} & \textbf{Months} & \textbf{Main activities} &
\textbf{Deliverable at the end of the phase}\\
\midrule
\endhead
\bottomrule
\endfoot
1 & M1--M4 & Port Awave-3D and GPUZIP to the HPC SDK (nvc++); re-integrate
nvCOMP/Bitcomp; reproduce the V100 baselines on Santos Dumont; establish
the block/\emph{tile} decomposition with a per-block precision tag;
per-block \emph{roofline} profiling (Nsight Compute) on the RTX 5090 and
Santos Dumont. & GPUZIP/Awave-3D ported and validated under the HPC SDK,
with numerical parity confirmed against the V100 baselines and the block
decomposition extended to carry precision \emph{tags}; per-block
\emph{roofline} profiling report.\\
\addlinespace
2 & M4--M10 & Instrument the solver to record FP dynamic-range statistics
(exponent histogram, minimum/maximum magnitude, variance) per block and
per time step from FP64 reference runs on the six datasets; correlate with
velocity-model attributes; build the \emph{a priori} precision-zoning
estimator; validate or refute the physical assumptions of the four
criteria. & FP dynamic-range characterization database (space-time, six
datasets) and \emph{a priori} precision-zoning estimator derived from the
velocity model, with a report on which model attributes are predictive and
which of the four criteria's assumptions hold empirically.\\
\addlinespace
3 & M10--M20 & Design, implement, and validate the four tagging criteria
(boundary/PML, energy, adjoint sensitivity, wavefront-trailing) informed
by Phase 2; compensated (Kahan) summation in the low-precision tiles;
precision-interface artifact testing; precision $\times$ compression
factorial study (four conditions) to define the unified per-tile error
budget. & Precision-AMR module with the four criteria implemented and
validated (PSNR, SSIM, and mean relative error vs. FP64 on the wavefield
and the migrated image), documented precision $\times$ compression
factorial methodology, and differential checkpointing integrated as the
storage mechanism for tagged tiles; open-source GPUZIP v3.0-alpha release
and reproducibility package (REDU).\\
\addlinespace
4 & M19--M27 & Apply the validated criteria on the RTX 5090; per-tile
\emph{pipeline} selecting between native tensor cores, mixed-precision GEMM
(cuBLAS), iterative refinement (cuSOLVER), and Ozaki-scheme FP64 emulation
in accuracy-critical tiles; quantify the performance $\times$ accuracy
trade-off against the uniform-Ozaki baseline; bound errors jointly
(compression $\times$ precision) across the six datasets and three
checkpointing algorithms. & Operational Precision-AMR \emph{pipeline} on
hardware without native FP64, with the selective-emulation gain quantified
against the uniform-Ozaki baseline and errors bounded vs. FP64 across the
six datasets.\\
\addlinespace
5 & M26--M31 & Incorporate into GPUZIP a compression library with native
AMD support (choice between hipCOMP-core, HPDR with MGARD/ZFP, or another,
decided at the start of the phase); port the tagging and \emph{tiling}
logic to the AMD Instinct MI300A nodes via ROCm/HIP; re-run the evaluation
protocol (criteria and precision $\times$ compression factorial) with the
AMD compressor; test direct transfer of the Phase 2 \emph{a priori}
estimator; compare performance $\times$ accuracy against the NVIDIA
baselines. & GPUZIP with a working AMD compression path (CDNA3) and a
Precision-AMR portability report for the MI300A, characterizing how much
of the gain is architecture-independent and how much is specific to the
CUDA/Ozaki/nvCOMP stack.\\
\addlinespace
6 & M31--M36 & Benchmark cuFFT-based pseudo-spectral operators as an
alternative update rule in tagged tiles; integrate the Precision-AMR
module (as an independent library), differential checkpointing, and
mixed-precision paths into a unified GPUZIP v3.0 release; multi-node
scaling evaluation ($\leq$8 GPUs) on Santos Dumont; final analysis,
write-up, and submission. & GPUZIP v3.0 release (MIT license, GitHub) with
the Precision-AMR library, the cuFFT operator \emph{benchmark}, and the
multi-node scaling evaluation; submitted paper and final public
reproducibility package (REDU).\\
\end{longtable}

\section{Description of the Activities Carried Out by the Team}

% TODO: section to be written last. Per the APR roteiro, include a brief
% description (up to one paragraph each) of the activities of:
%  - Principal Investigator (Sandro Rigo);
%  - Associated Researchers (Herve Yviquel);
%  - requested scholarship holders (if any);
%  - students and collaborators not funded by FAPESP (Thiago Maltempi).

\emph{[Placeholder --- section to be written last. It will describe, in up
to one paragraph per person, the activities of the Principal Investigator,
the Associated Researchers, the requested scholarship holders, and the
other collaborators not funded by FAPESP.]}

% ============================================================================
% OLD BLOCK (pre-FAPESP restructuring) --- kept for reference.
% Remove only after final approval of the restructured version.
% ============================================================================
\iffalse

\section*{Abstract}

Scientific simulations such as Reverse Time Migration (RTM) solve partial
differential equations (PDEs) demanding FP64/FP32 accuracy, yet modern GPU
architectures increasingly optimize for low-precision AI workloads (FP8/BF16)
and expose new memory hierarchies, tensor cores, and mixed-precision math
libraries. Prior work on RTM acceleration has reduced floating-point precision
uniformly across the whole computational domain; separately, adaptive mesh
refinement (AMR) has been explored for seismic wave modeling to vary grid
spacing according to local wavelength. This project proposes a third, largely
unexplored axis of adaptivity: Precision-AMR (\pamr{}), a spatially
block-structured scheme that varies floating-point precision --- not mesh
resolution --- according to physically motivated, per-block error-tolerance
criteria (proximity to absorbing boundaries, wavefield energy density,
adjoint-sensitivity magnitude, and a time-dependent band trailing the
propagating wavefront).

Unlike mesh-based AMR, where numerical artifacts are localized at coarse-fine
grid interfaces, precision-induced round-off noise propagates with the wave
itself, so \pamr{}'s tagging criteria must target where injected error will not
reach the imaging condition or the FWI gradient, and must control round-off
accumulation over the thousands of time steps typical of RTM/FWI. We build
\pamr{} on our prior GPUZIP framework, which accelerated RTM on V100 systems
via GPU-side lossy compression (nvCOMP/Bitcomp) and checkpoint prefetching,
cutting host--GPU transfer volume by up to 80$\times$ and delivering up to
5.1$\times$ end-to-end speedup. Differential checkpointing, cuBLAS/cuSOLVER
mixed-precision linear algebra, Ozaki-scheme FP64 emulation, and cuFFT-based
spectral operators are retained as supporting techniques, now applied
selectively, per tile, according to the criteria that \pamr{} establishes,
rather than uniformly across the domain.

The 12-month programme combines two computational resources. An RTX PRO 6000
Blackwell GPU, requested as dedicated local hardware, supports day-to-day
development and the selective-emulation study central to Phase 3, since its
negligible native FP64 throughput makes the precision-zoning question
unavoidable rather than optional. The Santos Dumont supercomputer
(LNCC/SINAPAD) --- whose partitions span NVIDIA Volta V100, Hopper H100, and
Grace-Hopper GH200 accelerators alongside AMD Instinct MI300A (CDNA3) APUs ---
supports the large-scale validation sweep across datasets, checkpointing
algorithms, and \pamr{} tagging criteria, and lets us test directly whether the
tagging methodology generalizes beyond a single GPU vendor. Outcomes include an
open-source GPUZIP v3.0 release built around a Precision-AMR module, and a
peer-reviewed publication.

\paragraph{Keywords:} reverse time migration; adaptive precision;
mixed-precision computing; floating-point arithmetic; adjoint PDE solvers;
block-structured adaptivity; cross-architecture portability; GPU computing;
lossy compression; checkpointing.

\fi
% ============================================================================
% END OF OLD BLOCK (Abstract)
% ============================================================================

% ============================================================================
% OLD BLOCK (pre-FAPESP restructuring) --- kept for reference.
% Content moved to the "Computational Platforms" subsection at the end of
% the "Scientific and Technological Challenges..." section (after Phase 6),
% since the FAPESP roteiro has no dedicated section for this.
% Remove only after final approval of the restructured version.
% ============================================================================
\iffalse

\section{Computational Platforms}

\paragraph{HPC SDK} (nvc++/nvfortran, CUDA-X math libraries, Nsight
Systems/Compute). Experience: team uses CUDA and Nsight Systems routinely.
Adoption: migrate the Awave-3D build from clang to nvc++ and validate numerical
parity against existing baselines. Nsight Compute provides kernel-level roofline
analysis, now applied per spatial tile, to identify which tiles are
memory-bandwidth-bound versus compute-bound and to objectively scope where a
lower-precision tile actually pays off.

\paragraph{RTX PRO 6000 Blackwell} (dedicated hardware, requested). This GPU's
5th-generation tensor cores with FP8/FP4 support and negligible native FP64
throughput make it the central platform for \pamr{}: because full-domain FP64
emulation would be expensive everywhere, the scientific question this project
answers is precisely which tiles need it and which do not. Its 96\,GB GDDR7 ECC
memory also supports local development; as dedicated, queue-free hardware it
enables fast iterative debugging that a shared batch system is not well suited
for.

\paragraph{Santos Dumont} (LNCC/SINAPAD, Tier-0 national supercomputer,
Petr\'opolis-RJ). Provides the large-scale validation environment. Its
BullSequana X1000 partition includes 94 nodes with 4$\times$ NVIDIA Volta V100
GPUs each --- matching the hardware used for the original GPUZIP baselines ---
plus a dedicated AI node with 8$\times$ V100 16GB over NVLink. Its newer
BullSequana XH3000 partition adds 62 nodes with 4$\times$ NVIDIA Hopper H100 SXM
80GB over NVLink, 36 NVIDIA Grace-Hopper GH200 superchip nodes, and 18 nodes
with 2$\times$ AMD Instinct MI300A APUs (CDNA3 GPU cores, 128\,GB HBM3 each).
This heterogeneity lets us generate FP64 ground truth and run the full \pamr{}
criteria sweep at scale on NVIDIA hardware, and separately test whether the
tagging methodology and block-tiling approach generalize to a non-NVIDIA
architecture.

\paragraph{nvCOMP / Bitcomp} (GPU-side lossy compressor). Experience: integrated
and validated in GPUZIP v2.0, where Bitcomp delivered the best
accuracy-compression balance on V100. Adoption here: the storage mechanism for
differential checkpointing, now indexed by the same tile identifiers used for
precision tagging, so each tile carries a single, jointly-tuned
precision-and-compression error budget.

\paragraph{cuBLAS (mixed-precision extensions) and cuSOLVER (iterative
refinement).} Both ship inside the HPC SDK. Experience: team has used
BLAS/LAPACK routines in prior HPC work. Adoption: applied only within tiles that
\pamr{} flags as containing dense sub-kernels needing accuracy beyond native low
precision (FP32 compute with FP64 accumulation, and cuSOLVER iterative
refinement).

\paragraph{Ozaki-scheme FP64 emulation} via low-precision tensor cores.
Adoption: reserved for the subset of tiles that \pamr{}'s sensitivity criteria
mark as accuracy-critical (near the wavefront, high adjoint-sensitivity
magnitude), rather than applied uniformly --- directly testing whether selective
emulation is cheaper than emulating everywhere while preserving image fidelity.
Recent work has demonstrated up to 2.3$\times$ speedup over native FP64 DGEMM on
the RTX PRO 6000 Blackwell~\cite{uchino2025ozaki}.

\paragraph{cuFFT.} Experience: used in prior spectral operator work. Adoption:
benchmarked as an alternative to FDM stencils for wavefield updates within
selected, precision-tagged tiles, characterising the memory-bandwidth versus
compute trade-off relative to the standard finite-difference path.

\paragraph{AMD Instinct MI300A} (CDNA3, Santos Dumont, exploratory). As an
additional validation activity, we assess how much of the \pamr{} framework ---
the block/tile decomposition, the four tagging criteria, and the joint
precision-compression error budget --- transfers to AMD's ROCm/HIP ecosystem
(rocBLAS, hipBLAS, and AMD's own reduced-precision matrix paths), without
assuming feature parity with the CUDA-specific Ozaki-scheme emulation or
nvCOMP/Bitcomp compression path. The goal is to determine whether Precision-AMR
is a portable methodology or an NVIDIA-specific optimization.

\fi
% ============================================================================
% END OF OLD BLOCK (Computational Platforms)
% ============================================================================

% ============================================================================
% UNNECESSARY SECTION FOR THE APR --- kept commented, not a restructuring
% "old block". "Dataset and Model" was required by a different FAPESP call
% (not the APR roteiro) and has no equivalent in the APR's mandatory
% structure. Kept here only for reference in case it is needed again; user
% decision on 2026-09-04.
% ============================================================================
\iffalse

\section{Dataset and Model}

\begin{itemize}
  \item Size: $\approx$X\,TB total (to be confirmed from REDU). Checkpoint pools
        reached $\approx$330\,GB host RAM on V100; a single Marmousi3D checkpoint
        is $\approx$500\,MB.
  \item Data type: seismic velocity models and traces (numerical grids).
  \item Source: Salt (SEG/EAGE) and Marmousi3D are open datasets; the Large
        dataset is in-house.
  \item Confidentiality: non-confidential, no personal data, no NDA.
  \item AI model parameters: not applicable (numerical PDE solver).
  \item Storage location: Universidade Estadual de Campinas (UNICAMP), Campinas,
        Brazil; staged for upload to Santos Dumont's parallel file systems for
        large-scale runs.
\end{itemize}

\fi
% ============================================================================
% END OF UNNECESSARY SECTION (Dataset and Model)
% ============================================================================

% ============================================================================
% OLD BLOCK (pre-FAPESP restructuring) --- kept for reference.
% Content already incorporated into the "Problem Statement" section above.
% Remove only after final approval of the restructured version.
% ============================================================================
\iffalse

\section{Introduction}

Hardware is trending toward low precision for AI, but PDE-based scientific
simulation requires high precision to remain trustworthy. RTM, a seismic imaging
method that solves the wave equation by finite differences over thousands of
time steps, is a canonical example: geophysicists use the migrated image to
guide expensive drilling decisions, so numerical fidelity is non-negotiable. RTM
is also memory-intensive, requiring hundreds of gigabytes of wavefield state and
relying on checkpointing to trade memory against recomputation in its
forward/backward adjoint structure.

Our prior work addressed the dominant bottleneck --- host--GPU checkpoint
transfer over PCIe --- with GPUZIP~\cite{maltempi2025checkpointing,gpuzip}, which
compresses checkpoints on the GPU and prefetches them proactively, reaching up to
5.1$\times$ speedup on V100 while preserving image quality (SSIM $\approx$ 1,
near-zero relative error). Separately, reduced-precision arithmetic has already
been shown to be viable for seismic modeling, RTM, and FWI:
Fabien-Ouellet~\cite{fabienouellet2020half} found that a stable half-precision
(FP16) implementation of the elastic wave equation is achievable through
appropriate scaling, that the resulting seismogram error is a negligible fraction
of total seismic energy and largely incoherent with seismic phases, and that this
noise does not adversely affect FWI or RTM image quality. However, that reduction
is applied uniformly across the entire domain. Gao and
Harms~\cite{gao2023half} identified a specific hazard of that approach: naive
half-precision time stepping accumulates round-off error through a disguised
recursive summation in the update rule, degrading solution quality, with
compensated (Kahan) summation as an effective remedy.

At the same time, adaptive mesh refinement (AMR) has been explored for seismic
wave simulation and, more recently, explicitly for FWI and RTM workflows, to
avoid the waste inherent in a uniform grid: dispersion and dissipation errors
require grid spacing proportional to the shortest wavelength, so a uniform mesh
is forced to use the finest spacing everywhere, even where the local velocity
field would tolerate a much coarser one. AMR addresses this by refining
$\Delta x$ locally. No comparable, physically-motivated scheme exists for
floating-point precision: existing mixed-precision work in seismic imaging is
either global and uniform (as above), or adaptive only at the level of an
individual linear-algebra operation (mixed-precision GEMM, iterative refinement,
Ozaki-scheme emulation), not tied to a spatial decomposition of the wavefield
itself.

This project proposes to close that gap with Precision-AMR (\pamr{}): an
adaptive-precision analogue of AMR that keeps mesh resolution fixed and instead
varies the floating-point format assigned to each block (tile) of the domain,
based on criteria that predict whether the round-off error introduced in that
tile can reach the imaging condition or the FWI gradient. This is a structurally
different problem from mesh-based AMR. In AMR, the characteristic artifact is a
localized reflection at the coarse-fine grid interface, because refinement
changes the discretization itself. In \pamr{}, no discretization boundary is
introduced; instead, a spatially varying level of round-off noise is injected
into a hyperbolic PDE, and that noise propagates together with the wave. Tagging
criteria therefore cannot rely solely on local smoothness (as AMR refinement
indicators do); they must instead approximate where an error, once injected, will
decay, be absorbed, or otherwise fail to reach the receivers or the sensitivity
kernel before the simulation ends.

We propose four candidate criteria, detailed in Methods: proximity to the
absorbing boundary (PML), where energy is already being discarded on purpose;
local wavefield energy density, which governs how much a given tile's error can
contribute to the final correlation; adjoint-sensitivity magnitude, reusing the
gradient computation that FWI already performs; and a time-dependent band
trailing the propagating wavefront, analogous to a moving shock detector in
astrophysical or combustion AMR codes. Because RTM/FWI already require a
block/tile decomposition of the domain for checkpointing, \pamr{} reuses that
same decomposition as its tagging granularity, unifying the precision budget and
the lossy-compression budget of differential checkpointing into a single, jointly
tuned per-tile accuracy target. Mixed-precision GEMM, cuSOLVER iterative
refinement, and Ozaki-scheme emulation remain essential, but are now applied
selectively, only where \pamr{}'s criteria say they are needed, which is the
natural way to exploit GPUs such as the RTX PRO 6000 Blackwell that offer
negligible native FP64 throughput.

Because the tagging criteria are motivated by the physics of wave propagation and
adjoint sensitivity rather than by any particular vendor's arithmetic units, we
treat their portability as an explicit question rather than an assumption. Access
to the Santos Dumont supercomputer --- whose partitions include AMD Instinct
MI300A accelerators alongside several NVIDIA generations --- lets us test that
question directly, rather than merely claim it.

\fi
% ============================================================================
% END OF OLD BLOCK (Introduction)
% ============================================================================

% ============================================================================
% OLD BLOCK (pre-FAPESP restructuring) --- kept for reference.
% Content reorganized (with new Phase 2 for empirical characterization, a
% dedicated Phase 5 for cross-architecture validation, and the timeline
% extended to 36 months) in the new "Scientific and Technological
% Challenges..." section above.
% Remove only after final approval of the restructured version.
% ============================================================================
\iffalse

\section{Methods}

\subsection{Phase 1 --- Infrastructure and Baseline (M1--M2)}

Port Awave-3D and GPUZIP to the HPC SDK (nvc++), re-integrate nvCOMP/Bitcomp, and
reproduce V100 baselines on Santos Dumont's V100 partition, matching the original
GPUZIP hardware.

\begin{itemize}
  \item Establish a block/tile decomposition of the domain, extending the
        existing checkpoint-block structure so that it can also carry a per-tile
        precision tag.
  \item Profile every kernel with Nsight Systems and Nsight Compute, on both the
        RTX PRO 6000 and Santos Dumont, confirming that FDM stencils are
        memory-bandwidth-bound at the tile granularity before any precision
        changes are attempted.
\end{itemize}

\subsection{Phase 2 --- Precision-AMR: Criteria, Tagging, and Validation (M2--M6)}

This is the core scientific phase, run at scale on Santos Dumont's H100
partition. We design, implement, and validate four candidate tagging criteria for
spatial precision zoning:

\begin{itemize}
  \item \textbf{Boundary criterion:} tiles inside or adjacent to the
        PML/absorbing layer are tagged for reduced precision by default, since
        their contribution is already being damped.
  \item \textbf{Energy criterion:} tiles are tagged by a threshold on local
        wavefield energy density, so that regions of low amplitude (before
        wavefront arrival, or after geometrical spreading has attenuated the
        signal) tolerate lower precision.
  \item \textbf{Sensitivity criterion (FWI):} tiles are tagged using the
        magnitude of the local adjoint-sensitivity kernel, reusing the gradient
        computation itself as the tagging signal, in direct analogy to a
        Berger--Oliger truncation-error estimator in classical AMR.
  \item \textbf{Wavefront-trailing criterion:} a time-dependent band immediately
        behind the propagating wavefront is tagged for reduced precision, while
        the leading edge --- where phase accuracy governs the imaging condition
        --- retains full precision.
\end{itemize}

For each criterion, we implement compensated (Kahan) summation in the
low-precision tiles' time-stepping update to control round-off accumulation over
the full simulation, following the mechanism identified by Gao and
Harms~\cite{gao2023half}. We explicitly test for a ``precision-interface''
artifact --- analogous to coarse-fine grid-interface reflections in mesh-based
AMR --- via residual analysis at tile boundaries between differing precisions.
Finally, each tile is assigned jointly a floating-point format and a checkpoint
compression policy (full snapshot versus compressed delta), with combined
accuracy measured by PSNR, SSIM, and mean relative error against the FP64
baseline, on both the raw wavefield and the final migrated image. This phase
subsumes differential checkpointing, which becomes the storage mechanism for
\pamr{}-tagged tiles rather than a separate deliverable.

\subsection{Phase 3 --- Precision-AMR without Native FP64 (M7--M10)}

The tagging criteria validated in Phase 2 are applied on the RTX PRO 6000,
hardware with negligible native FP64 throughput. Per tile, the pipeline selects
between native low-precision tensor-core execution, cuBLAS mixed-precision GEMM
(FP32 compute with FP64 accumulation) or cuSOLVER iterative refinement for dense
sub-kernels, and Ozaki-scheme FP64 emulation for tiles flagged as
accuracy-critical.

\begin{itemize}
  \item This directly operationalises the motivation for exploiting AI-first GPUs
        that lack native FP64: expensive emulation is concentrated only where
        \pamr{}'s criteria say it is needed.
  \item We quantify the speedup/accuracy trade-off relative to a uniform-Ozaki
        baseline (emulating FP64 everywhere) to determine how much selective
        tiling actually saves.
  \item All errors are bounded jointly (compression $\times$ precision
        interaction) against the FP64 baseline using PSNR, SSIM, and mean
        relative error, across the same three datasets and three checkpointing
        algorithms (Revolve, zCut, Uniform) used throughout the project.
  \item Exploratory cross-architecture activity: the same tagging and tiling logic
        is ported to Santos Dumont's AMD Instinct MI300A nodes via ROCm/HIP, to
        test whether the criteria and error-budget framework hold on a CDNA3
        architecture without native access to Ozaki-scheme emulation or
        nvCOMP/Bitcomp.
\end{itemize}

\subsection{Phase 4 --- Spectral Operators, Integration, and Release (M10--M12)}

\begin{itemize}
  \item Benchmark cuFFT-based pseudo-spectral derivative operators as an
        alternative update rule for selected, precision-tagged tiles, as a
        secondary technique complementing the core \pamr{} contribution.
  \item Integrate the Precision-AMR module, differential checkpointing, and
        mixed-precision paths into a unified GPUZIP v3.0 release.
  \item Multi-node scaling evaluation on Santos Dumont ($\leq$8 GPUs).
  \item Final analysis, write-up, and submission, with Precision-AMR as the
        central scientific claim and differential checkpointing, mixed-precision
        linear algebra, and spectral operators presented as supporting
        techniques.
\end{itemize}

\fi
% ============================================================================
% END OF OLD BLOCK (Methods)
% ============================================================================

\section{Evaluation and Dissemination}

This section describes how the project's results will be evaluated and
disseminated. Evaluation covers the datasets and the experimental
protocol; dissemination covers the channels for scientific publication,
open-source software, and outreach.

\subsection{Datasets}

The experiments use six datasets. Five are open synthetic models, widely
used as benchmarks by the seismic imaging community, and the sixth is a
large in-house model, used mainly in the scalability tests.

The term \emph{SEG/EAGE open data}~\cite{segopendata} denotes the body of
synthetic velocity models and seismic datasets released publicly by the
\emph{Society of Exploration Geophysicists} (SEG) and the \emph{European
Association of Geoscientists and Engineers} (EAGE) for benchmarking
processing and imaging algorithms (migration, velocity analysis, multiple
suppression), for the study of wave propagation in complex geological
media, and for comparing hardware platforms. These data are
non-confidential and free for research use.

\begin{itemize}
  \item \textbf{Marmousi3D:} a three-dimensional extension of the Marmousi
        model~\cite{versteeg1994marmousi}, a synthetic model of complex
        geological structure, with intense faulting and strong lateral
        velocity variation, which has become a standard benchmark for FWI
        and RTM.
  \item \textbf{SEG/EAGE Salt:} a 3D model with a large, complex-geometry
        salt body over a faulted sedimentary sequence, part of the
        SEG/EAGE 3-D Modeling Series~\cite{aminzadeh1997segeage}. It is
        the canonical subsalt imaging benchmark.
  \item \textbf{SEG/EAGE Overthrust:} a 3D model of a thrust belt
        (compressional tectonics), with steeply dipping structure and
        strong lateral velocity variation, also from the SEG/EAGE 3-D
        Modeling Series~\cite{aminzadeh1997segeage}. It tests imaging
        under strong shallow heterogeneity.
  \item \textbf{Sigsbee 2A and Sigsbee 2B:} 2D synthetic subsalt models,
        inspired by the Sigsbee Escarpment in the Gulf of Mexico and
        released by the SMAART JV consortium~\cite{paffenholz2002sigsbee}.
        The 2A model is a constant-density acoustic model without
        free-surface multiples; the 2B model adds strong internal and
        free-surface multiples. They are classic benchmarks for subsalt
        imaging under low signal-to-noise conditions.
  \item \textbf{Large:} a large in-house model (velocities and traces on
        numerical grids), used mainly in the multi-node scalability tests,
        as it requires the largest volume of wavefield state and
        checkpoints.
\end{itemize}

All are numerical grids of seismic velocity models and traces, with no
personal data and no confidentiality restriction.

\subsection{Evaluation protocol}

All experiments span the six datasets above, three checkpointing
algorithms (Revolve, zCut, Uniform), and the four \pamr{} tagging criteria
(individually and in combination). Kernel-level profiling uses Nsight
Compute (roofline, memory-traffic breakdowns); timeline profiling uses
Nsight Systems. Implementation is in C/C++ with CUDA under the HPC SDK, with OpenMP
Cluster for multi-node scaling. Accuracy is evaluated on both the wavefield and
the final RTM image, since a tiling that preserves wavefield PSNR/SSIM may still,
in principle, shift the migrated image if error is spatially correlated with
reflectors. For the cross-architecture activity, the same accuracy metrics are
computed on Santos Dumont's AMD MI300A partition and compared against the NVIDIA
baselines, to characterise how much of the achieved speedup and accuracy
trade-off is architecture-independent.

\subsection{Dissemination strategy}

The project's results will be disseminated through three complementary
channels:

\begin{itemize}
  \item \textbf{Scientific publications:} papers in relevant conferences
        and journals in high-performance computing and seismic imaging,
        with Precision-AMR as the central contribution. The venues under
        consideration are detailed in the Expected Results section.
  \item \textbf{Open-source software:} the Precision-AMR library and GPUZIP
        v3.0 will be released into the public domain on GitHub, or an
        equivalent platform, under the MIT license, together with the
        reproducibility package (datasets, Nsight profiles, quality
        metrics) made available via UNICAMP's REDU.
  \item \textbf{Outreach:} the Laboratory of Computer Systems (LSC), to
        which the Principal Investigator belongs, regularly publicizes
        research results on LinkedIn; this channel will be used to broaden
        the reach of the publications and of the software release.
\end{itemize}

% ============================================================================
% OLD BLOCK (pre-FAPESP restructuring) --- kept for reference. The
% "Projected outcomes" paragraph has already been incorporated into the new
% "Expected Results" section above (right after "Problem Statement"). The
% timeline (month/phase table) will be reused in the future "Execution
% Timeline" section; the "Dependencies" paragraph has not yet been
% relocated. Remove only after final approval of the restructured version.
% ============================================================================
\iffalse

\section{Expected Results}

Timeline (12 months):

\begin{longtable}{@{}p{0.12\textwidth}p{0.5\textwidth}p{0.19\textwidth}p{0.09\textwidth}@{}}
\toprule
\textbf{Month} & \textbf{Activity} & \textbf{Hardware} & \textbf{Phase}\\
\midrule
\endfirsthead
\toprule
\textbf{Month} & \textbf{Activity} & \textbf{Hardware} & \textbf{Phase}\\
\midrule
\endhead
\bottomrule
\endfoot
M1--M2 & Port Awave-3D + GPUZIP to HPC SDK (nvc++); reproduce V100 baselines;
establish block/tile decomposition; Nsight Compute roofline profiling per tile.
& RTX PRO 6000; Santos Dumont (V100) & 1\\
\addlinespace
M2--M5 & Design and implement Precision-AMR tagging criteria (PML proximity,
energy density, adjoint sensitivity, wavefront-trailing band); compensated
summation; precision-interface artifact testing. & Santos Dumont (H100) & 2\\
\addlinespace
M5--M6 & Unify precision and compression error budget via differential
checkpointing; multi-node scaling; open-source GPUZIP v3.0-alpha release;
reproducibility package (REDU). & Santos Dumont (H100, multi-node) & 2\\
\addlinespace
M7--M10 & Apply Precision-AMR tiling without native FP64: selective cuBLAS
mixed-precision GEMM, cuSOLVER iterative refinement, and Ozaki-scheme FP64
emulation only in accuracy-critical tiles. Exploratory port of the tagging
framework to AMD Instinct MI300A via ROCm/HIP. & RTX PRO 6000; Santos Dumont
(AMD MI300A, exploratory) & 3\\
\addlinespace
M10--M12 & cuFFT spectral-operator benchmark on selected tiles; unified GPUZIP
v3.0 release; analysis, write-up, and paper submission. & RTX PRO 6000; Santos
Dumont & 4\\
\end{longtable}

Projected outcomes: an open-source GPUZIP v3.0 release (MIT license, GitHub)
built around a Precision-AMR module, with differential checkpointing,
mixed-precision GEMM/Ozaki-scheme paths, and cuFFT spectral operators integrated
as supporting components; a peer-reviewed publication with Precision-AMR as the
central contribution (target: SBAC-PAD / Euro-Par / IJHPCA); and a public
reproducibility package (datasets, Nsight profiles, quality metrics) via UNICAMP
REDU.

Dependencies: this hardware request is not conditional on being combined with
another grant. Large-scale validation runs on the Santos Dumont supercomputer
(LNCC/SINAPAD) through LNCC's own computational allocation program (SINAPAD),
independent of this request. Existing institutional support funds personnel; none
of these sources duplicates the RTX PRO 6000 request.

\fi
% ============================================================================
% END OF OLD BLOCK (Expected Results)
% ============================================================================

\section{Project Support Details}

We request a dedicated local server equipped with an NVIDIA RTX 5090 GPU
(32\,GB GDDR7) and 64\,GB of RAM, with an estimated cost of up to
R\$~55,000 in Brazil, within the per-item equipment ceiling set by
FAPESP's Regular Research Grant rules. This GPU shares with same-generation
professional Blackwell cards the same memory bandwidth (1,792\,GB/s) and
the same 5th-generation tensor cores, which makes it representative for
kernel development on a memory-bandwidth-bound solver like ours, at a
fraction of the cost of an equivalent professional card. Its negligible
native FP64 throughput makes it the scientifically relevant target for the
selective-emulation study central to Phase 4: it forces the question of
whether Precision-AMR's selective use of Ozaki-scheme emulation can match
full-domain emulation at a fraction of the cost. As dedicated, queue-free
hardware, it also supports the iterative day-to-day development that a
shared batch system is not well suited for. Large-scale validation, that
is, the \pamr{} criteria sweep across datasets and checkpointing
algorithms, FP64 baseline generation, multi-GPU scaling, and the AMD
MI300A portability test, runs separately on the Santos Dumont
supercomputer (LNCC/SINAPAD), accessed through LNCC's own allocation
program and outside the scope of this request.

\subsection*{Resource Request}

\begin{longtable}{@{}p{0.3\textwidth}p{0.62\textwidth}@{}}
\toprule
\textbf{Resource} & \textbf{Request}\\
\midrule
\endhead
\bottomrule
\endfoot
Physical hardware & 1$\times$ server with an NVIDIA RTX 5090 GPU (32\,GB
GDDR7) and 64\,GB of RAM, dedicated local development hardware, used
throughout the project (estimated cost: up to R\$~55,000)\\
\end{longtable}

Compute time on Santos Dumont is requested separately through LNCC's own project
allocation program (SINAPAD) and is not part of this hardware request.

\bibliographystyle{plainnat}
\bibliography{mixed-precision-amr-en}

\end{document}
